\documentclass[11pt,a4paper]{article}

\usepackage[margin=1in]{geometry}
\usepackage[T1]{fontenc}
\usepackage[utf8]{inputenc}
\usepackage{lmodern}
\usepackage{xcolor}
\usepackage{booktabs}
\usepackage{longtable}
\usepackage{array}
\usepackage{enumitem}
\usepackage{amsmath}
\usepackage{amssymb}
\usepackage{listings}
\usepackage[hidelinks,breaklinks=true]{hyperref}

\definecolor{codebg}{RGB}{248,248,248}
\definecolor{codeframe}{RGB}{220,220,220}
\definecolor{commentgreen}{RGB}{0,110,55}
\definecolor{keywordblue}{RGB}{0,70,150}
\definecolor{stringred}{RGB}{150,35,35}

\lstdefinestyle{cppstyle}{
  language=C++,
  basicstyle=\ttfamily\scriptsize,
  keywordstyle=\color{keywordblue}\bfseries,
  commentstyle=\color{commentgreen},
  stringstyle=\color{stringred},
  numbers=left,
  numberstyle=\tiny\color{gray},
  stepnumber=1,
  numbersep=6pt,
  breaklines=true,
  breakatwhitespace=false,
  tabsize=2,
  keepspaces=true,
  showstringspaces=false,
  backgroundcolor=\color{codebg},
  frame=single,
  rulecolor=\color{codeframe},
  columns=fullflexible,
  captionpos=b
}

\newcommand{\codepath}[1]{\path{#1}}

\title{Clarification Report on PBC Parameters, Hashing, Full Keys, and State Information}
\author{VANET ns-3 Security Research Workspace}
\date{April 29, 2026}

\begin{document}
\maketitle

\begin{abstract}
This short report answers four implementation-level questions about the current
pairing-based cryptography (PBC) security flow in the VANET ns-3 module. It
focuses only on the main points needed for explanation: how the KGC and TA
secrets are selected, how the hash function \(H_1\) is implemented, whether the
vehicle secret \(x_i\) and group elements are in the expected algebraic sets, and
how \(Q_i\) differs from the state information \(\theta\) used during signature
generation.
\end{abstract}

\section{Quick Answer Summary}

\begin{longtable}{p{0.24\linewidth}p{0.68\linewidth}}
\toprule
\textbf{Question} & \textbf{Short implementation answer} \\
\midrule
Stage 1: are \(s\) and \(t\) random or from \(\mathbb{Z}_q^*\)? &
They are selected randomly using PBC \codepath{element_random} after being
initialized in \codepath{Zr}. In this implementation, \codepath{Zr} is the scalar
field of order \(r\) from the Type A pairing parameters. The code does not
hard-code 160 bits; the approximately 160-bit scalar order comes from the
pairing parameter \codepath{r}. \\
\midrule
Partial key: what is \(H_1\)? Is it SHA-256? &
For partial key generation, \(H_1\) is implemented as \codepath{HashToG1}, which
uses PBC \codepath{element_from_hash} to map bytes into \(G_1\). SHA-256 is used
separately for PID2 masking, not for the \(\Pi_i=H_1(PID_i\parallel P_{pub})\)
step. \\
\midrule
Full key: is \(x_i\in\mathbb{Z}_q^*\) and are \(P,S_i\in G_1\)? &
Yes in the practical PBC notation: \(x_i\) is initialized in \codepath{Zr} and
randomly sampled; \(P\), \(S_i\), \(Q_i\), \(W_i\), \(X\), \(R_i\), and
\(\psi_i\) are initialized as \codepath{G1} elements. Strictly speaking, the code
samples from \(\mathbb{Z}_r\), not explicitly \(\mathbb{Z}_r^*\), because it does
not resample if zero occurs. \\
\midrule
How is \(Q\) chosen as state info? &
\(Q_i\) is not the state information. \(Q_i=x_iP\) is the vehicle public
verification key. The state information is \(\theta\), implemented as the
\codepath{StateInfo} string with default value \codepath{traffic-alert}. The
signature uses both \(Q_i\) and \(\theta\). \\
\bottomrule
\end{longtable}

\section{Question 1: Selection of \texorpdfstring{\(s\)}{s} and \texorpdfstring{\(t\)}{t}}

\subsection{Formula}

In the five-stage PBC protocol, the KGC selects a master secret \(s\), and the TA
selects a master secret \(t\):

\[
  s \xleftarrow{\$} \mathbb{Z}_r, \qquad P_{pub}=sP
\]

\[
  t \xleftarrow{\$} \mathbb{Z}_r, \qquad T_{pub}=tP
\]

Some papers write the scalar set as \(\mathbb{Z}_q^*\). The current code uses the
PBC library's scalar field name \codepath{Zr}. In the Type A parameters, \(r\) is
the group/scalar order. Therefore, for this implementation, the closest mapping
is:

\[
  \mathbb{Z}_q^*\text{ in paper notation} \approx \mathbb{Z}_r^*\text{ in PBC notation}
\]

The code samples randomly from \(\mathbb{Z}_r\). It does not explicitly check for
non-zero, so it is not a strict \(\mathbb{Z}_r^*\) sampler. The probability of
sampling zero is negligible, about \(1/r\), but a strict implementation can add a
resampling check.

\subsection{Code Evidence}

The generator point \(P\) is initialized in \(G_1\), and then \(s\) and \(t\) are
initialized in \codepath{Zr} and randomized.

\begin{lstlisting}[style=cppstyle,caption={Setup code for P, s, Ppub, t, and Tpub.}]
element_init_G1(m_P, m_pairing);
const char kGeneratorSeed[] = "VANET_PBC_GENERATOR_SEED_V1";
element_from_hash(m_P,
                  const_cast<void*>(static_cast<const void*>(kGeneratorSeed)),
                  static_cast<int>(sizeof(kGeneratorSeed) - 1));
m_hasP = true;

// KGC setup: s in Zr and Ppub = sP
element_init_Zr(m_s, m_pairing);
element_random(m_s);
element_init_G1(m_Ppub, m_pairing);
element_mul_zn(m_Ppub, m_P, m_s);

// TA setup: t in Zr and Tpub = tP
element_init_Zr(m_t, m_pairing);
element_random(m_t);
element_init_G1(m_Tpub, m_pairing);
element_mul_zn(m_Tpub, m_P, m_t);
\end{lstlisting}

The Type A pairing parameters used by LTE/NR include the scalar order \(r\):

\begin{lstlisting}[style=cppstyle,caption={Type A pairing parameters used by the scenario.}]
std::string pbcParams = R"(type a
q 5016916071032417911863923687783915420518698834988272735086782217043042258214706291487114949230487855429610676959493682417064495802878387205742192057730447
h 6865426548812704696509566582849235473375093961437097302961574578811789461292366439277723796099293963270768
r 730750818665452757176057050065048642452048576511
exp2 159
exp1 110
sign1 1
sign0 -1
)";
\end{lstlisting}

\subsection{Conclusion}

\begin{itemize}[leftmargin=*]
  \item \(s\) and \(t\) are selected randomly.
  \item They are selected from PBC \codepath{Zr}, whose order is defined by the pairing parameter \(r\).
  \item The implementation does not manually choose a 160-bit prime in code; the 160-bit scalar order comes from the pairing parameters.
  \item In paper notation this may be written as \(\mathbb{Z}_q^*\); in this code the corresponding PBC scalar field is \(\mathbb{Z}_r\), with \(r\) as the group order.
\end{itemize}

\section{Question 2: What Algorithm Implements \texorpdfstring{\(H_1\)}{H1}?}

\subsection{Formula}

Partial private key generation uses:

\[
  PID_i = PID1_i \parallel PID2_i \parallel T_i
\]

\[
  \Pi_i = H_1(PID_i \parallel P_{pub})
\]

\[
  S_i = s\Pi_i
\]

Here \(H_1\) must output a point in \(G_1\), not just a fixed-size byte digest.
That means \(H_1\) is a hash-to-group function.

\subsection{Code Evidence}

The current implementation first encodes the pseudonym, then appends
\(P_{pub}\), and finally maps the resulting bytes into \(G_1\).

\begin{lstlisting}[style=cppstyle,caption={PID encoding and partial key generation.}]
bool
PbcCrypto::EncodePid(const std::vector<uint8_t>& pid1Bytes,
                     const std::vector<uint8_t>& pid2Bytes,
                     uint64_t tiUs,
                     std::vector<uint8_t>& pidBytes)
{
  pidBytes.clear();
  WriteUint16(pidBytes, static_cast<uint16_t>(pid1Bytes.size()));
  WriteBytes(pidBytes, pid1Bytes);
  WriteUint16(pidBytes, static_cast<uint16_t>(pid2Bytes.size()));
  WriteBytes(pidBytes, pid2Bytes);
  WriteUint64(pidBytes, tiUs);
  return true;
}

bool
PbcCrypto::ComputePi(const std::vector<uint8_t>& pidBytes,
                     std::vector<uint8_t>& piBytes)
{
  std::vector<uint8_t> ppubBytes = GetPpubBytes();
  std::vector<uint8_t> input;
  input.insert(input.end(), pidBytes.begin(), pidBytes.end());
  input.insert(input.end(), ppubBytes.begin(), ppubBytes.end());
  return HashToG1(input, piBytes);
}

bool
PbcCrypto::ComputePartialKey(const std::vector<uint8_t>& pid1Bytes,
                             const std::vector<uint8_t>& pid2Bytes,
                             uint64_t tiUs,
                             std::vector<uint8_t>& siBytes)
{
  std::vector<uint8_t> pidBytes;
  EncodePid(pid1Bytes, pid2Bytes, tiUs, pidBytes);

  std::vector<uint8_t> piBytes;
  ComputePi(pidBytes, piBytes);

  element_t pi;
  element_t si;
  element_init_G1(pi, m_pairing);
  element_init_G1(si, m_pairing);
  element_from_bytes(pi, piBytes.data());
  element_mul_zn(si, pi, m_s);   // Si = s * Pi
  return EncodeElement(si, siBytes);
}
\end{lstlisting}

The actual \(H_1\) implementation is:

\begin{lstlisting}[style=cppstyle,caption={HashToG1 implementation used for H1.}]
bool
PbcCrypto::HashToG1(const std::vector<uint8_t>& data, std::vector<uint8_t>& out)
{
  if (!m_initialized || data.empty())
  {
    return false;
  }
  element_t tmp;
  element_init_G1(tmp, m_pairing);
  element_from_hash(
      tmp, const_cast<void*>(static_cast<const void*>(data.data())),
      static_cast<int>(data.size()));
  bool ok = EncodeElement(tmp, out);
  element_clear(tmp);
  return ok;
}
\end{lstlisting}

\subsection{Important Distinction from SHA-256}

SHA-256 is used in PID2 generation, not in \(H_1\) for partial key generation.
The PID2 masking formula is:

\[
  q_i=tPID1_i,
  \qquad h_i=\operatorname{SHA256}(q_i\parallel T_{pub}),
  \qquad PID2_i=ID_i\oplus h_i[0..3]
\]

\begin{lstlisting}[style=cppstyle,caption={SHA-256 is used for PID2 masking.}]
element_mul_zn(q, pid1, m_t);
EncodeElement(q, qBytes);
EncodeElement(m_Tpub, tpubBytes);

std::vector<uint8_t> hashInput;
hashInput.insert(hashInput.end(), qBytes.begin(), qBytes.end());
hashInput.insert(hashInput.end(), tpubBytes.begin(), tpubBytes.end());
std::vector<uint8_t> digest = EccCrypto::Sha256(hashInput.data(), hashInput.size());

pid2Bytes[0] = idBytes[0] ^ digest[0];
pid2Bytes[1] = idBytes[1] ^ digest[1];
pid2Bytes[2] = idBytes[2] ^ digest[2];
pid2Bytes[3] = idBytes[3] ^ digest[3];
\end{lstlisting}

\subsection{Conclusion}

\begin{itemize}[leftmargin=*]
  \item \(H_1\) is not an encryption algorithm.
  \item \(H_1\) is implemented as \codepath{HashToG1} using PBC \codepath{element_from_hash}.
  \item The input to \(H_1\) for partial key generation is \(PID_i\parallel P_{pub}\).
  \item SHA-256 is used for PID2 masking, not for \(\Pi_i=H_1(PID_i\parallel P_{pub})\).
\end{itemize}

\section{Question 3: Full Key Generation and Algebraic Sets}

\subsection{Formula}

The vehicle full key generation should satisfy:

\[
  x_i \xleftarrow{\$} \mathbb{Z}_r, \qquad Q_i=x_iP
\]

For signing, the implementation should use:

\[
  S_i,P,Q_i,W_i,X,R_i,\psi_i \in G_1
\]

\[
  x_i,w_i \in \mathbb{Z}_r
\]

\[
  \psi_i = S_i + x_iX + w_iR_i
\]

\subsection{Code Evidence: Vehicle Verification Key}

\begin{lstlisting}[style=cppstyle,caption={Vehicle xi is sampled in Zr and Qi is computed in G1.}]
bool
PbcCrypto::GenerateVerificationKey(std::vector<uint8_t>& xiBytes,
                                   std::vector<uint8_t>& qiBytes)
{
  element_t xi;
  element_t qi;
  element_init_Zr(xi, m_pairing);  // xi in Zr
  element_init_G1(qi, m_pairing);  // Qi in G1
  element_random(xi);
  element_mul_zn(qi, m_P, xi);     // Qi = xi * P

  bool ok = EncodeElement(xi, xiBytes) && EncodeElement(qi, qiBytes);
  element_clear(qi);
  element_clear(xi);
  return ok;
}
\end{lstlisting}

The vehicle calls this after it receives the KGC partial key \(S_i\):

\begin{lstlisting}[style=cppstyle,caption={Vehicle generates Qi after receiving partial key.}]
if (siLen > 0)
{
  m_partialKey = siBytes;
  if (m_enablePbc)
  {
    if (!m_pbc.GenerateVerificationKey(m_secretXi, m_publicQi))
    {
      m_secretXi.clear();
      m_publicQi.clear();
    }
  }
}
\end{lstlisting}

\subsection{Code Evidence: Signature Elements}

\begin{lstlisting}[style=cppstyle,caption={Signature generation uses Zr scalars and G1 group elements.}]
element_t si;
element_t xi;
element_t qi;
element_t wi;
element_t wiPoint;
element_t x;
element_t ri;
element_t xiX;
element_t wiRi;
element_t psi;

element_init_G1(si, m_pairing);       // Si in G1
element_init_Zr(xi, m_pairing);       // xi in Zr
element_init_G1(qi, m_pairing);       // Qi in G1
element_init_Zr(wi, m_pairing);       // wi in Zr
element_init_G1(wiPoint, m_pairing);  // Wi in G1
element_init_G1(x, m_pairing);        // X in G1
element_init_G1(ri, m_pairing);       // Ri in G1
element_init_G1(xiX, m_pairing);      // xi * X in G1
element_init_G1(wiRi, m_pairing);     // wi * Ri in G1
element_init_G1(psi, m_pairing);      // psi in G1

element_from_bytes(si, siBytes.data());
element_from_bytes(xi, xiBytes.data());
element_mul_zn(qi, m_P, xi);

element_random(wi);
element_mul_zn(wiPoint, m_P, wi);     // Wi = wi * P

ComputeX(thetaBytes, xBytes);
ComputeRi(thetaBytes, messageBytes, pidBytes, qiBytes, wiBytes, riBytes);
element_from_bytes(x, xBytes.data());
element_from_bytes(ri, riBytes.data());

element_mul_zn(xiX, x, xi);           // xi * X
element_mul_zn(wiRi, ri, wi);         // wi * Ri
element_set(psi, si);
element_add(psi, psi, xiX);
element_add(psi, psi, wiRi);          // psi = Si + xiX + wiRi
\end{lstlisting}

\subsection{Conclusion}

\begin{itemize}[leftmargin=*]
  \item Yes, \(x_i\) is currently generated as a random \codepath{Zr} element.
  \item Yes, \(P\), \(S_i\), \(Q_i\), \(W_i\), \(X\), \(R_i\), and \(\psi_i\) are handled as \codepath{G1} elements.
  \item The only strict-theory caveat is that the code does not explicitly reject zero for \(x_i\) or \(w_i\). For strict \(\mathbb{Z}_r^*\), a resampling loop can be added.
\end{itemize}

\section{Question 4: Is \texorpdfstring{\(Q_i\)}{Qi} the State Information?}

\subsection{Clarification}

\(Q_i\) is not the state information. The implementation has two different
objects:

\begin{itemize}[leftmargin=*]
  \item \(Q_i=x_iP\): the vehicle's public verification key.
  \item \(\theta\): shared state information used by the PBC signature formula.
\end{itemize}

The signature formula uses both:

\[
  X = H_1(\theta)
\]

\[
  R_i = H_1(\theta \parallel M_i \parallel PID_i \parallel Q_i \parallel W_i)
\]

\[
  \psi_i = S_i + x_iX + w_iR_i
\]

So \(Q_i\) is included inside the hash for \(R_i\), but \(Q_i\) is not equal to
\(\theta\).

\subsection{Code Evidence: State Info Attribute}

\begin{lstlisting}[style=cppstyle,caption={Vehicle state info theta is a string attribute.}]
.AddAttribute("StateInfo",
              "Shared state information theta for PBC V2V signatures.",
              StringValue("traffic-alert"),
              MakeStringAccessor(&VehicleApp::m_stateInfo),
              MakeStringChecker())
\end{lstlisting}

\subsection{Code Evidence: Signature Input Uses Both Theta and Qi}

\begin{lstlisting}[style=cppstyle,caption={Vehicle signs using theta and sends Qi separately.}]
std::vector<uint8_t> pidBytes;
std::vector<uint8_t> thetaBytes(m_stateInfo.begin(), m_stateInfo.end());
std::vector<uint8_t> wiBytes;
std::vector<uint8_t> psiBytes;

m_pbc.EncodePid(m_pid1, m_pid2, m_pidTiUs, pidBytes);
m_pbc.SignMessage(m_partialKey,
                  m_secretXi,
                  pidBytes,
                  thetaBytes,
                  payload,
                  wiBytes,
                  psiBytes);

VanetPbcAuthHeader auth;
auth.SetState(thetaBytes);   // theta
auth.SetPid1(m_pid1);
auth.SetPid2(m_pid2);
auth.SetTimestampUs(m_pidTiUs);
auth.SetQi(m_publicQi);      // Qi = xi * P
auth.SetWi(wiBytes);
auth.SetPsi(psiBytes);
packet->AddHeader(auth);
\end{lstlisting}

The relay batches packets by the state \(\theta\), not by \(Q_i\):

\begin{lstlisting}[style=cppstyle,caption={Relay batches by state theta and verifies with Qi.}]
std::string batchKey(auth.GetState().begin(), auth.GetState().end());
auto& batch = m_batches[batchKey];

entry.stateBytes = auth.GetState();
entry.pidBytes = pidBytes;
entry.qiBytes = auth.GetQi();
entry.wiBytes = auth.GetWi();
entry.psiBytes = auth.GetPsi();

verified = m_pbc.AggregatePsi(psiList, aggregatePsiBytes) &&
           m_pbc.VerifyAggregate(pidList,
                                 qiList,
                                 wiList,
                                 messageList,
                                 batch.entries.front().stateBytes,
                                 aggregatePsiBytes);
\end{lstlisting}

\subsection{Conclusion}

\begin{itemize}[leftmargin=*]
  \item \(Q_i\) is chosen by sampling \(x_i\in\mathbb{Z}_r\) and computing \(Q_i=x_iP\).
  \item State information is \(\theta\), implemented as \codepath{m_stateInfo} and defaulting to \codepath{traffic-alert}.
  \item The relay groups signatures by \(\theta\) because aggregate verification uses one common \(X=H_1(\theta)\) for the batch.
  \item \(Q_i\) is sent separately in the PBC authentication header and is used in \(R_i\) and in the aggregate verification equation.
\end{itemize}

\section{Optional Strictness Improvement}

For strict mathematical compliance with \(\mathbb{Z}_r^*\), the code can resample
until a non-zero scalar is obtained. Conceptually:

\begin{lstlisting}[style=cppstyle,caption={Optional strict non-zero scalar sampling pattern.}]
element_init_Zr(m_s, m_pairing);
do
{
  element_random(m_s);
}
while (element_is0(m_s));
\end{lstlisting}

The same pattern can be applied to \(t\), \(d_i\), \(x_i\), and \(w_i\). The
current implementation is still practically safe because the probability of
zero is approximately \(1/r\), but the resampling check makes it exactly match
\(\mathbb{Z}_r^*\) notation.

\end{document}
